\begin{proofbox}

\textbf{\textcolor{CProof}{思路提示}}

利用旋转体体积公式 $V=\pi\int_{-b}^{b}x^2\,dy$，其中 $x$ 为椭圆上对应 $y$ 处的 $x$ 坐标。由椭圆方程解出 $x^2$ 并积分。

\medskip
\textbf{\textcolor{CProof}{正式推导}}
\begin{description}[style=nextline, leftmargin=2.8em, labelsep=0.5em, itemsep=0.55em]
    \item[\textbf{步骤一（建立积分表达式）：}] 
    将旋转体视为垂直 $y$ 轴的圆盘叠加，圆盘半径 $r(y)=|x|$，体积微元 $dV=\pi x^2\,dy$。沿 $y$ 从 $-b$ 到 $b$ 积分：
    \[
    V = \pi \int_{-b}^{b} x^2 \,dy
    \]
    \par\textbf{理由：} 旋转体由平面曲线绕轴旋转生成，可用微元法。

    \item[\textbf{步骤二（写出 $x^2$ 表达式）：}]
    由椭圆方程 $\dfrac{x^2}{a^2}+\dfrac{y^2}{b^2}=1$，解出 $x^2$：
    \[
    x^2 = a^2\left(1-\frac{y^2}{b^2}\right)
    \]
    \par\textbf{理由：} $a>0$，两边同乘 $a^2$ 得。

    \item[\textbf{步骤三（代入并积分）：}]
    将 $x^2$ 代入积分式：
    \[
    V = \pi \int_{-b}^{b} a^2\left(1-\frac{y^2}{b^2}\right) dy
    \]
    提出常数 $a^2$：
    \[
    V = \pi a^2 \int_{-b}^{b} \left(1-\frac{y^2}{b^2}\right) dy
    \]
    \par\textbf{理由：} $a^2$ 为常数，可提出积分号。

    \item[\textbf{步骤四（计算定积分）：}]
    先求不定积分，再代入上下限：
    \[
    \begin{aligned}
    \int_{-b}^{b} \left(1-\frac{y^2}{b^2}\right) dy
    &= \left[ y - \frac{y^3}{3b^2} \right]_{-b}^{b} \\
    &= (b - \frac{b}{3}) - (-b + \frac{b}{3}) \\
    &= 2\left(b - \frac{b}{3}\right) \\
    &= \frac{4}{3}b
    \end{aligned}
    \]
    \par\textbf{理由：} 牛顿-莱布尼茨公式，注意 $(-b)^3=-b^3$。

    \item[\textbf{步骤五（得体积公式）：}]
    将积分结果代入：
    \[
    \begin{aligned}
    V &= \pi a^2 \cdot \frac{4}{3}b \\
    &= \frac{4}{3}\pi a^2 b
    \end{aligned}
    \]
    \par\textbf{理由：} 化简乘积得到最终体积公式。
\end{description}

\medskip


\medskip
\textbf{\textcolor{CProof}{新增方法二：沿 $x$ 轴切成椭圆薄片（异向截面积分）}}

\textbf{思路：} 原证明是垂直于旋转轴 $y$ 切圆片；现在改成垂直于 $x$ 轴切片。在旋转椭球中，$x=x_0$ 截到的不是圆，而是位于 $yz$ 平面的椭圆。面积能算出来，也就能累积出体积。

\begin{description}[style=nextline, leftmargin=2.8em, labelsep=0.5em, itemsep=0.55em]
  \item[\textbf{步骤一（三维方程）：}]
  平面椭圆绕 $y$ 轴旋转时，原来的 $x^2$ 对应三维空间中到 $y$ 轴的距离平方 $x^2+z^2$，因此实体为
  \[
    E=\left\{(x,y,z):\frac{x^2+z^2}{a^2}+\frac{y^2}{b^2}\le1\right\}.
  \]
  \par\textbf{理由：} 绕 $y$ 轴旋转不改变 $y$ 坐标，距 $y$ 轴的距离是 $\sqrt{x^2+z^2}$。

  \item[\textbf{步骤二（确定 $x=x_0$ 的截面）：}]
  固定 $x=x_0$（$-a<x_0<a$），移项得
  \[
     \frac{y^2}{b^2}+\frac{z^2}{a^2}
     \le 1-\frac{x_0^2}{a^2}.
  \]
  记 $q=\sqrt{1-\frac{x_0^2}{a^2}}$，即
  \[
     \frac{y^2}{(bq)^2}+\frac{z^2}{(aq)^2}\le1.
  \]
  \par\textbf{理由：} 这是 $yz$ 平面内的椭圆区域；其 $y$ 方向半轴为 $bq$，$z$ 方向半轴为 $aq$。

  \item[\textbf{步骤三（计算截面积）：}]
  根据椭圆面积 $\pi\times\text{长半轴}\times\text{短半轴}$（两个半轴谁长均可），有
  \[
     A(x_0)=\pi (bq)(aq)
        =\pi ab\left(1-\frac{x_0^2}{a^2}\right).
  \]
  \par\textbf{理由：} 此处面积是\textbf{椭圆面积} $\pi pq$，不是 $\pi r^2$。

  \item[\textbf{步骤四（沿 $x$ 方向积分）：}]
  \[
  \begin{aligned}
    V_y&=\int_{-a}^{a}A(x)\,dx
      =\pi ab\int_{-a}^{a}\left(1-\frac{x^2}{a^2}\right)\,dx\\
    &=\pi ab\left[x-\frac{x^3}{3a^2}\right]_{-a}^{a}\\
    &=\pi ab\left(\frac{2a}{3}+\frac{2a}{3}\right)
      =\boxed{\frac43\pi a^2b}.
  \end{aligned}
  \]
  \par\textbf{理由：} 积分区间是 $x\in[-a,a]$；两种不同方向切片给出相同体积。
\end{description}

\textbf{考试速写（沿 $x$ 轴切）：}
\[
   \boxed{A(x)=\pi ab\left(1-\frac{x^2}{a^2}\right)}
   \quad\Longrightarrow\quad
   V_y=\pi ab\int_{-a}^{a}\left(1-\frac{x^2}{a^2}\right)\,dx
      =\frac43\pi a^2b.
\]

\medskip
\textbf{\textcolor{CProof}{新增方法三：由球体体积伸缩得到（最少步骤）}}

\textbf{几何思路：} 取一个半径为 $b$ 的球 $B$，其体积 $\frac43\pi b^3$ 已知。将球的 $x$、$z$ 坐标各乘 $a/b$，$y$ 坐标不变；椭球就出现了。面积伸缩为 $\left(\frac ab\right)^2$，长度 $y$ 不变，所以整个体积也乘这个倍数。

\begin{description}[style=nextline, leftmargin=2.8em, labelsep=0.5em, itemsep=0.55em]
  \item[\textbf{步骤一（从球出发）：}]
  \[
      B:\quad u^2+v^2+w^2\le b^2,
      \qquad V_B=\frac43\pi b^3.
  \]

  \item[\textbf{步骤二（作空间伸缩）：}]
  定义
  \[
      x=\frac ab\,u,\qquad y=v,\qquad z=\frac ab\,w.
  \]
  反代 $u=\frac ba x,\ v=y,\ w=\frac ba z$ 可得
  \[
      \frac{x^2+z^2}{a^2}+\frac{y^2}{b^2}\le1.
  \]
  \par\textbf{理由：} 恰是绕 $y$ 轴旋转的椭球实体；注意\textbf{要同时改变 $x,z$ 两方向}。

  \item[\textbf{步骤三（体积伸缩）：}]
  空间线性伸缩把体积乘以三个方向伸缩比的乘积
  \[
   \begin{aligned}
      V_y&=V_B\cdot\frac ab\cdot1\cdot\frac ab
        =\frac43\pi b^3\cdot\frac{a^2}{b^2}\\
        &=\boxed{\frac43\pi a^2b}.
   \end{aligned}
  \]
  \par\textbf{理由：} 直观上可先在每个 $y$ 高度把圆截面积放大 $\left(\frac ab\right)^2$，再沿 $y$ 方向叠加；严格表述是线性变换的体积比例 $\lvert\det\mathrm{diag}(a/b,1,a/b)\rvert=(a/b)^2$。
\end{description}

\begin{center}
  \resizebox{0.95\linewidth}{!}{% C026_sphere_stretch.tikz.tex -- 直接 \input，需先 \usepackage{tikz}。
% 图示为 z=0 的二维母线截面。三维中 x 和 z 必须同步伸缩。
\begin{tikzpicture}[x=0.79cm,y=0.79cm,>=stealth,every node/.style={font=\small}]
  \begin{scope}[shift={(0,0)}]
    \fill[blue!3,rounded corners=8pt] (-3.3,-2.55) rectangle (3.3,2.8);
    \draw[blue!20,rounded corners=8pt] (-3.3,-2.55) rectangle (3.3,2.8);
    \draw[->,gray!70] (-2.8,0)--(2.85,0) node[below,font=\scriptsize] {$x$};
    \draw[->,gray!70] (0,-2.2)--(0,2.42) node[left,font=\scriptsize] {$y$};
    \fill[blue!15] (0,0) circle [radius=1.7];
    \draw[blue!75!black,thick] (0,0) circle [radius=1.7];
    \draw[blue!75!black,line width=1.1pt,->] (0,0)--(1.7,0);
    \draw[blue!75!black,line width=1.1pt,->] (0,0)--(0,1.7);
    \node[font=\scriptsize,blue!75!black] at (1.04,-0.35) {$b$};
    \node[font=\scriptsize,blue!75!black] at (-0.23,1.35) {$b$};
    \node[font=\scriptsize] at (0,-2.22) {半径 $b$ 的球（母线是圆）};
  \end{scope}
  \draw[->,line width=1.5pt,purple!70!black] (3.65,0.45)--(6.4,0.45)
       node[midway,above,align=center,font=\scriptsize] {$x,z$ 均乘 $a/b$};
  \begin{scope}[shift={(10.15,0)}]
    \fill[teal!3,rounded corners=8pt] (-3.3,-2.55) rectangle (3.3,2.8);
    \draw[teal!20,rounded corners=8pt] (-3.3,-2.55) rectangle (3.3,2.8);
    \draw[->,gray!70] (-2.95,0)--(2.98,0) node[below,font=\scriptsize] {$x$};
    \draw[->,gray!70] (0,-2.2)--(0,2.42) node[left,font=\scriptsize] {$y$};
    \fill[teal!13] (0,0) ellipse [x radius=2.54,y radius=1.7];
    \draw[teal!75!black,thick] (0,0) ellipse [x radius=2.54,y radius=1.7];
    \draw[teal!75!black,line width=1.1pt,->] (0,0)--(2.54,0);
    \draw[teal!75!black,line width=1.1pt,->] (0,0)--(0,1.7);
    \node[font=\scriptsize,teal!70!black] at (1.44,-0.35) {$a$};
    \node[font=\scriptsize,teal!70!black] at (-0.23,1.35) {$b$};
    \node[font=\scriptsize] at (0,-2.22) {绕 $y$ 轴旋转椭球（母线是椭圆）};
  \end{scope}
  \node[font=\scriptsize,purple!70!black] at (5.05,-2.25) {体积乘 $(a/b)^2$};
\end{tikzpicture}
}
\end{center}

\textbf{最快记忆／复原法：}
\[
   \boxed{\text{球 }\frac43\pi R^3\ \longrightarrow\
          \text{椭球 }\frac43\pi(\text{半轴}_1)(\text{半轴}_2)(\text{半轴}_3)}
    =\frac43\pi\cdot a\cdot b\cdot a.
\]
这里的推广是\textbf{椭球体积公式}；此题是其中两个半轴相等的旋转椭球，不是任意椭球都能由某个平面椭圆绕轴旋转得到。

\textbf{\textcolor{CProof}{结论回扣}}

\textbf{我们通过积分严格证明了椭圆绕 $y$ 轴旋转所得旋转椭球体积为 $V_y = \frac{4}{3}\pi a^2 b$。}

\end{proofbox}
